\section{Open theoretical questions}
\label{sec:open_questions}

This section discusses three fundamental open questions in TGP Path~B
that require theoretical rather than numerical resolution.

\subsection{Q1: Yukawa source mechanism}
\label{sec:yukawa_source}

\paragraph{Question.}
Why does a body of mass $m$ act as a source of the Yukawa scalar field
$\phi$ with coupling $C = m/(2\sqrt{\pi}\,m_{\rm Pl})$?
In TGP Path~B, this is taken as an axiom.  No derivation from first
principles has been given.

\paragraph{Scalar-field formulation (standard}.
In a standard massive scalar field theory with Lagrangian
$\mathcal{L} = -\frac{1}{2}(\partial_\mu\phi)^2 - \frac{1}{2}m_{\rm sp}^2\phi^2 - g\phi T$,
the coupling to the trace of the energy-momentum tensor $T=-\rho$ (for
non-relativistic matter) gives the field equation
$(\Box - m_{\rm sp}^2)\phi = g\rho$,
with the Yukawa solution $\phi = -\frac{g}{4\pi}\frac{m e^{-m_{\rm sp}r}}{r}$
for a point mass $m$.

The two-body potential from scalar exchange is then
$V_{12} = g^2\phi_1(r_{12})m_2 = -\frac{g^2 m_1 m_2}{4\pi}\frac{e^{-m_{\rm sp}r}}{r}$.
Matching to TGP: $g^2/(4\pi) = C_1 C_2 / (m_1 m_2) = 1/(4\pi m_{\rm Pl}^2)$,
giving $g = 1/m_{\rm Pl}$ — this is the Planck-scale coupling.

\paragraph{Three-body vertex from $\phi^3$ term.}
The three-body TGP interaction arises if the scalar Lagrangian contains
a cubic self-interaction $\mathcal{L}_3 = -\lambda\phi^3/6$.  The one-loop
diagram with one $\phi^3$ vertex and three external matter lines generates
the Feynman integral $I_Y$ computed in the preceding sections.  The
coupling scales as $C^3 = (g m/m_{\rm Pl})^3 = (m/m_{\rm Pl})^3/(4\pi)^{3/2}$.

\paragraph{Identification with Dilaton or Brans-Dicke scalar.}
A natural candidate is the dilaton field $\phi$ of string theory or the
scalar partner in Brans-Dicke gravity.  In Brans-Dicke theory, the scalar
$\chi$ couples to matter as $g_{\rm BD}\chi T$, with $g_{\rm BD}\sim1/m_{\rm Pl}$.
The Yukawa mass $m_{\rm sp}$ would arise from a non-minimal scalar potential
$V(\chi)\ni\frac{1}{2}m_{\rm sp}^2\chi^2$.

\paragraph{Status.}
The coupling structure of TGP Path~B is consistent with a standard
dilaton-like scalar with coupling $g=1/m_{\rm Pl}$ and mass $m_{\rm sp}$.
The derivation from a specific string or modified-gravity Lagrangian
is an open problem.  The key observation is that the TGP 2-body
potential IS the Newtonian gravitational potential at $r\ll1/m_{\rm sp}$,
so the scalar-exchange mechanism is a completion of Newtonian gravity
at sub-Planck scales.

\subsection{Q2: Dark matter and rotation curves}
\label{sec:dark_matter}

\paragraph{Question.}
Can TGP with $m_{\rm sp}\sim1/{\rm kpc}$ explain galactic rotation curves
without dark matter?

\paragraph{Answer: No, for ordinary matter.}
The TGP two-body coupling for ordinary baryonic matter is
$C = m/(2\sqrt{\pi}\,m_{\rm Pl}) \sim 10^{-19}$--$10^{-20}$
(for masses $m\sim10^{-27}$--$10^{-26}$~kg).
The two-body TGP potential is:
\begin{equation}
  V_{\rm TGP}(r) = -\frac{G m_1 m_2}{r}\,e^{-m_{\rm sp}r}.
\end{equation}
At $r\ll1/m_{\rm sp}$ (sub-Yukawa range) this is Newtonian gravity.
At $r\gg1/m_{\rm sp}$ the potential is \emph{exponentially suppressed},
giving \emph{less} gravity than Newton.  Flat rotation curves require
\emph{more} gravity at large $r$ (not less), so Yukawa screening
cannot explain dark-matter rotation curves for any $m_{\rm sp}$.

The three-body correction scales as $C^3\sim10^{-60}$ and is
completely negligible for galactic dynamics.

\paragraph{TGP dark matter scenario.}
A TGP dark matter candidate would require particles with
$C\sim1$ (near-Planck-mass objects, $m\sim m_{\rm Pl}$) in sufficient
density to contribute to galactic rotation.  Such particles would:
\begin{itemize}
\item Interact gravitationally (via the same Yukawa) with baryons
  with $C\sim10^{-19}$, giving effectively no observable coupling
  at galactic scales.
\item Be gravitationally self-interacting with $C^2\sim1$, potentially
  forming bound clusters (the Efimov-type states of this work).
\item Have no electroweak or strong interactions, consistent with
  collisionless dark matter at galaxy scales.
\end{itemize}

\paragraph{MOND connection.}
Modified Newtonian dynamics (MOND, Milgrom 1983) postulates a
transition at acceleration $a_0\approx1.2\times10^{-10}\,{\rm m/s}^2$.
TGP provides no derivation of $a_0$.  The TGP three-body force
gives a correction $\delta a/a\sim C\sim10^{-19}$ for baryonic matter,
far below $a_0/g_{\rm Newt}\sim10^{-4}$ at galactic edges.
TGP and MOND are presently independent frameworks; a connection would
require a theoretical bridge identifying $a_0$ with some combination
of $C$ and $m_{\rm sp}$.

\paragraph{Numerical comparison (ex31).}
A rotation curve computation (see ex31) confirms that Yukawa-screened
gravity (TGP) gives:
\begin{itemize}
\item Identical to Newtonian for $r\ll1/m_{\rm sp}$ (as expected).
\item Below Newtonian for $r\gtrsim1/m_{\rm sp}$ — opposite to observation.
\item CDM (NFW halo) and MOND both produce flat curves and fit
  observations; TGP alone does not.
\end{itemize}

\subsection{Q3: Path C — quantum defects and the Yukawa derivation}
\label{sec:path_c}

\paragraph{Question.}
TGP Path~C proposes that the Yukawa coupling of Path~B can be
\emph{derived} from the quantum theory of topological defects of the
scalar field $\phi$.  What is the mathematical structure of this
derivation?

\paragraph{Topological defect review (ex10-ex11).}
The TGP scalar field satisfies a nonlinear equation with defect solutions
(ex10-ex11).  The linearised tail of the defect is not Yukawa but
oscillatory:
\begin{equation}
  \phi(r)\sim\frac{A\sin(r)}{r} + \frac{B\cos(r)}{r}\quad
  (r\to\infty).
\end{equation}
This oscillatory tail arises from the Helmholtz-type linearisation
$(\nabla^2+1)\delta=0$ near $\phi=1$.  A Yukawa tail
($e^{-m_{\rm sp}r}/r$) would require a \emph{massive} linearised
equation $(\nabla^2-m_{\rm sp}^2)\delta=0$, which is not present
in the current TGP defect ODE.

\paragraph{Sketch of Path C.}
Path~C would require modifying the scalar Lagrangian (or adding
a mass term) so that the defect tail becomes Yukawa.  One approach:
\begin{enumerate}
\item Add a symmetry-breaking potential $V(\phi)=\frac{1}{2}m_{\rm sp}^2(\phi-1)^2$
  to the TGP Lagrangian.
\item The linearised equation becomes $(\nabla^2-m_{\rm sp}^2)\delta=0$
  near $\phi=1$, giving Yukawa decay.
\item The defect core acts as a Yukawa source with coupling $C$
  determined by the defect profile — this would \emph{derive}
  the TGP Path~B axiom.
\end{enumerate}
This sketch requires:
(a) A consistent modified TGP Lagrangian with both defect and Yukawa solutions,
(b) A derivation of $C(m_{\rm sp})$ from the defect quantisation,
(c) Consistency with the oscillatory solutions found in ex11
  (which would emerge in the limit $m_{\rm sp}\to0$).

\paragraph{Status.}
Path~C is mathematically open.  The key obstacle identified in ex11 is
that the \emph{standard} TGP defect ODE produces oscillatory (not Yukawa)
tails.  A modified Lagrangian is required.  The construction of a
consistent TGP Lagrangian connecting defect solutions to Yukawa couplings
is the primary open theoretical problem.

\subsection{Q4: Relation to general relativity}
\label{sec:gr_relation}

\paragraph{Question.}
TGP Path~B is a scalar field theory; GR is a tensor theory.
How do they coexist, and what are the observational signatures
distinguishing them?

\paragraph{Coexistence.}
In the standard picture, two independent fields can mediate gravity:
\begin{itemize}
\item \textbf{Tensor (graviton):} mediates spin-2 gravity, responsible
  for GR effects (light deflection, perihelion precession, gravitational waves).
\item \textbf{Scalar ($\phi$, TGP):} mediates additional spin-0 attraction,
  responsible for the Yukawa correction $e^{-m_{\rm sp}r}$.
\end{itemize}
The post-Newtonian parameter framework (PPN) accommodates both:
\begin{equation}
  \Phi(r) = -\frac{Gm}{r}\left(1 + \alpha_{\rm TGP}\,e^{-m_{\rm sp}r}\right),
\end{equation}
where $\alpha_{\rm TGP}=C^2/(Gm/r)\sim1$ is the ratio of TGP to Newtonian
coupling.  For TGP Path~B with $C=m/(2\sqrt{\pi}m_{\rm Pl})$, one has
$\alpha_{\rm TGP}\sim1$ at $r\ll1/m_{\rm sp}$, so TGP contributes
$\sim 100\%$ of Newtonian gravity at short range.

\paragraph{Observational distinction from GR.}
Key signatures of the scalar sector:
\begin{enumerate}
\item \textbf{Fifth-force violation of equivalence principle:}
  Scalar coupling $g=1/m_{\rm Pl}$ is universal (proportional to mass),
  so equivalence principle is preserved and TGP is equivalent to Newtonian
  for $r\ll1/m_{\rm sp}$.
\item \textbf{Yukawa deviation from $1/r^2$ force:}
  At $r\gtrsim1/m_{\rm sp}$, TGP predicts sub-Newtonian gravity.
  Laboratory and satellite tests of the inverse-square law constrain
  $m_{\rm sp}\gtrsim10^{-3}\,{\rm m}^{-1}$ ($\lambda_{\rm sp}\lesssim10^3\,{\rm m}$).
  Planck-scale $m_{\rm sp}\sim1/l_{\rm Pl}$ is far above this bound.
\item \textbf{Scalar gravitational waves:}
  The TGP scalar $\phi$ can emit monopole/dipole radiation in addition
  to GR's quadrupole.  The scalar wave energy is
  $P_{\rm scalar}\sim G^2 m^2 \ddot{m}^2_{\rm sp}/c^5$ — negligible
  for non-relativistic sources.
\item \textbf{Brans-Dicke parameter:}
  If TGP is identified with Brans-Dicke theory with parameter
  $\omega_{\rm BD}$, solar system tests require $\omega_{\rm BD}>40{,}000$.
  This constrains the scalar coupling; TGP with $g=1/m_{\rm Pl}$
  and $m_{\rm sp}\gg H_0$ (screening at sub-solar-system scales)
  is consistent with $\omega_{\rm BD}\to\infty$ (no long-range scalar modification).
\end{enumerate}

\paragraph{Summary.}
TGP Path~B is consistent with GR and all solar-system tests if
$m_{\rm sp}\gtrsim1/(1\,{\rm AU})\approx7\times10^{-12}\,{\rm m}^{-1}$.
For Planck-scale $m_{\rm sp}\sim l_{\rm Pl}^{-1}=6\times10^{34}\,{\rm m}^{-1}$,
the Yukawa range is $\lambda_{\rm sp}\sim10^{-35}$~m — far below any
observable scale except the Planck-mass N-body physics studied in this work.
TGP and GR coexist without contradiction at all observationally accessible scales.

\subsection{Summary of open questions}

\begin{table}[h]
\centering
\caption{Status of open theoretical questions in TGP Path~B.}
\label{tab:open_status}
\begin{tabular}{p{0.7cm}p{4cm}p{3.5cm}p{4cm}}
\hline
Q & Question & Status & Key obstacle \\
\hline
Q1 & Yukawa source derivation & Open & No Lagrangian for TGP \\
Q2 & Dark matter/rotation curves & Answered (negative) & TGP reduces gravity at large $r$; 3B negligible \\
Q3 & Path C: defect $\to$ Yukawa & Open & Oscillatory defect tail; need modified Lagrangian \\
Q4 & TGP vs GR & Addressed & Coexist; TGP is sub-Planck modification of gravity \\
Q5 & Full 2D quantum Efimov & Open & Equilateral approx.\ is upper bound; need 2D FD solver \\
\hline
\end{tabular}
\end{table}
